\section*{第 \ref{ch:b3:field-lab-methods-3} 章——野外与实验方法 III}
% Values from the Book 5 scratch calculator calc/ch31.py.

\begin{solution}{exo:b3:field-lab-methods-3:1}
(a)~$\qty{0.01}{\milli\gal} = \qty{e-7}{m.s^{-2}}$，且
$\num{e-7}/\num{9.81} = \num{1.0e-8}$：即一亿分之一。

(b)~$0.01/0.3086 = \qty{0.032}{m}$，约为 \qty{3.2}{cm}。

(c)~两种效应共同使重力改变
\qty{0.1967}{\milli\gal.m^{-1}}，即差值 $0.3086 - 0.1119$，因此
高度须改变 $0.01/0.1967 = \qty{5.1}{cm}$。
\end{solution}

\begin{solution}{exo:b3:field-lab-methods-3:2}
$\Delta m = \tfrac43\pi(\qty{50}{m})^3\times\qty{500}{kg.m^{-3}} =
\qty{2.62e8}{kg}$。峰值为 $G\Delta m/z^2 =
\num{6.674e-11}\times\num{2.62e8}/\num{e4} = \qty{1.75e-6}{m.s^{-2}}$，
即 \qty{0.175}{\milli\gal}。半宽为 $x_{1/2} = z/1.305 =
\qty{77}{m}$。
\end{solution}

\begin{solution}{exo:b3:field-lab-methods-3:3}
(a)~设 $AM = MN = NB = a$，$K = 2\pi/(1/a - 1/2a - 1/2a + 1/a) =
2\pi a = \qty{62.83}{m}$，且 $\rho_a = 62.83\times0.239/0.1 =
\qty{150}{\ohm\metre}$。(b)~$K = \pi(15^2 - 5^2)/(2\times5) = \pi\times
200/10 = \qty{62.83}{m}$。此时电极位于 $0$、$10$、$20$
及 \qty{30}{m}，按 AMNB 顺序排列，这是间距为
\qty{10}{m} 的温纳\omterm{def:b3:field-lab-methods-3:resistivity}{电极排列}。
\end{solution}

\begin{solution}{exo:b3:field-lab-methods-3:4}
(a)~$V = 0.2998/\sqrt5 = \qty{0.134}{m.ns^{-1}}$。(b)~$t =
2\times10/0.134 = \qty{149}{ns}$。(c)~$\delta = 0.0053\sqrt5/\num{e-3} =
\qty{11.9}{m}$，有效深度约为 \qty{36}{m}。(d)~在黏土中
$\delta = 0.0053\sqrt{20}/0.01 = \qty{2.4}{m}$，有效深度约为
\qty{7}{m}，且速度仅为 \qty{0.067}{m.ns^{-1}}。此处
近似在此处仍然成立：在 \qty{100}{MHz} 时，
$\omega\varepsilon = 2\pi\times\num{e8}\times20\times\num{8.85e-12} =
\qty{0.11}{S.m^{-1}}$，为 $\sigma$ 的十倍。
\end{solution}

\begin{solution}{exo:b3:field-lab-methods-3:5}
(a)~圆柱体的线密度为 $\lambda = \pi R^2\Delta\rho$，且
在其外部，引力与位于其轴上的线质量相同（由高斯
定理可知，共轴圆柱）：$\Delta g = 2G\pi R^2\Delta\rho\,z/(x^2
+ z^2)$。峰值一半处满足 $x^2 + z^2 = 2z^2$，故 $z =
x_{1/2}$。(b)~球体位于 $1.305\times100 = \qty{130}{m}$ 处，
圆柱体位于 \qty{100}{m} 处：同一剖面若按球体解释，则会把场源
估计得深 \qty{30}{\percent}。球体异常随距离衰减更快，
因此球体须埋藏更深，才能产生同样的宽度。
\end{solution}

\begin{solution}{exo:b3:field-lab-methods-3:6}
(a)~每个数据点代表 $\qty{100}{m}\times\qty{100}{m} =
\qty{e4}{m^2}$，故 $\iint\Delta g\,\mathrm{d}A \approx
\num{-438e-5}\times\num{e4} = \qty{-43.8}{m^3.s^{-2}}$，且 $\Delta m
= -43.8/(2\pi\times\num{6.674e-11}) = \qty{-1.04e11}{kg}$。(b)~$1.04/1.74
= \qty{60}{\percent}$。其余部分位于图幅之外异常的
两翼中：花岗岩深度超过一千米，且其
深部异常随距离衰减缓慢。扩大图幅
至足够远，或借助模型，即可恢复之。
\end{solution}

\begin{solution}{exo:b3:field-lab-methods-3:7}
(a)~对于薄岩墙，异常在 $x = \pm z$ 处降至峰值的一半：
其半峰全宽为 $2z = \qty{60}{m}$，跨越它的三个读数
所需间距至多为 \qty{20}{m}。(b)~$2\times
220/3 = \qty{147}{m}$。(c)~所需测点间距与感兴趣的最浅地质体的
深度成正比：按深部岩墙设计的测量，
间距为 \qty{150}{m}，将混淆浅部岩墙的异常，后者
宽度仅 \qty{60}{m}；按浅部地质体设计的测量
也能看到深部岩墙，但读数远多于所需。
\end{solution}

\begin{solution}{exo:b3:field-lab-methods-3:8}
(a)~$F = 80/8 = 10$，而 $\phi = F^{-1/m} = 10^{-1/2} = \num{0.316}$。
(b)~$10^{-1/1.8} = \num{0.278}$，$10^{-1/2.2} = \num{0.351}$。$m$ 改变
\num{0.2}，会使孔隙度约变化四个百分点；
电阻率的测量精度约为百分之几，对孔隙度的影响小得多。胶结指数
控制着不确定度，且须在同种岩石的样品上
进行标定。
\end{solution}

\begin{solution}{exo:b3:field-lab-methods-3:9}
(a)~当 $s = 2h_1$ 时，$(2nh_1/s)^2 = n^2$，故各项为 $k^n(1 +
n^2)^{-3/2}$：$0.3536k$、$0.0894k^2$、$0.0316k^3$。取 $k = 0.5$：
$\num{0.1768}$、$\num{0.0224}$、$\num{0.0040}$，且 $\rho_a/\rho_1 \approx
1 + 2\times0.2031 = \num{1.406}$（完整求和给出 \num{1.409}）。取 $k
= -0.5$ 时，各项符号交替，$-0.1768 + 0.0224 - 0.0040$，且
$\rho_a/\rho_1 \approx \num{0.683}$（完整求和 \num{0.685}）。(b)~当 $s \ll
h_1$ 时各项均为零，$\rho_a \to \rho_1$；当 $s \gg h_1$ 时
各括号项趋于~1，且 $1 + 2\sum k^n = 1 + 2k/(1 - k) = (1 + k)/(1 -
k) = \rho_2/\rho_1$。(c)~$\rho_2/\rho_1 = (1 + k)/(1 - k) = 3$ 与 $1/3$。
\end{solution}

\begin{solution}{exo:b3:field-lab-methods-3:10}
(a)~按 AMNB 顺序，间距为 $a$：$AM = a$，$BM = 2a$，$AN = 2a$，
$BN = a$，故括号项为 $2/a - 1/a = 1/a$，且 $K = 2\pi a$。
(b)~设 A、B 位于 $\mp s$，M、N 位于 $\mp l$：$AM = BN = s - l$，
$BM = AN = s + l$，括号项为 $2/(s - l) - 2/(s + l) = 4l/(s^2 -
l^2)$，且 $K = \pi(s^2 - l^2)/(2l)$。当 $l \ll s$ 时，$\Delta V_{MN}
\approx 2l\,E$，其中 $E = 2(-\mathrm{d}V/\mathrm{d}r)$ 为中心处的
两供电电极共同产生的场；于是 $\rho_a = K\Delta V/I \to
\pi s^2/(2l)\times 4l(-\mathrm{d}V/\mathrm{d}r)/I = 2\pi s^2(-
\mathrm{d}V/\mathrm{d}r)/I$。(c)~近似因子与精确因子之比
为 $s^2/(s^2 - l^2) = 1/(1 - 0.01) = 1.010$：即百分之一。
\end{solution}

\begin{solution}{exo:b3:field-lab-methods-3:11}
(a)~读数：\qty{0.010}{\milli\gal}。高度项：
$0.1967\times0.05 = \qty{0.010}{\milli\gal}$。密度项：布格
校正为 $0.04191\,h\,d$，$d$ 的单位为 \unit{g.cm^{-3}}，故
$0.04191\times50\times0.1 = \qty{0.210}{\milli\gal}$。地形项：
\qty{0.200}{\milli\gal}。合计：$\sqrt{0.010^2 + 0.010^2 + 0.210^2 +
0.200^2} = \qty{0.29}{\milli\gal}$。(b)~密度项与地形项
起主导作用，各约为读数误差的二十倍。欲探测
\qty{0.1}{\milli\gal} 的异常，板片的密度须在
样品上测定，或通过数据拟合确定（取最少反映
地形的值），并由精细高程模型做\omterm{def:b3:field-lab-methods-3:terrain}{地形校正}。
(c)~$\sqrt{0.010^2 + 0.010^2 + 0.210^2} = \qty{0.21}{\milli\gal}$。
\end{solution}

\begin{solution}{exo:b3:field-lab-methods-3:12}
(a)~与中心读数相符的质量为 $\Delta m = \Delta g(0)z^2/G
= \num{-4.16e-5}z^2/\num{6.674e-11}$：在 \qty{300}{m} 处为 \qty{-5.61e10}{kg}，
在 \qty{400}{m} 处为 \qty{-9.97e10}{kg}，
在 \qty{500}{m} 处为 \qty{-1.56e11}{kg}。预测读数在 $x = \pm300$ 和 $\pm600$ 处分别为
\num{-1.47} 和 \num{-0.37}、\num{-2.13} 和 \num{-0.71}、\num{-2.62}
和 \qty{-1.09}{\milli\gal}。\omterm{def:b3:field-lab-methods-3:inversion-practice}{拟合差}分别为 $\chi^2 = 1238$、$3.7$ 和
$856$。(b)~只有 \qty{400}{m} 处的模型在误差范围内拟合，其 $\chi^2 = 3.7$，
对应五个数据。(c)~可以：任何质量与中心相同的球对称体给出完全相同的异常，
其他形状的体也可能落在误差范围内。
在样品上测得的密度、钻孔或其他方法，
可以作出判定。
\end{solution}

\begin{solution}{pb:b3:field-lab-methods-3:1}
\textbf{1.} $(2846.708 - 2846.512)/\qty{2.0}{h} = \qty{0.098}{\milli\gal.h^{-1}}$。

\textbf{2.} $2851.075 - 0.098\times0.5 - 2846.512 = \qty{4.514}{\milli\gal}$；
$2850.497 - 0.088 - 2846.512 = \qty{3.897}{\milli\gal}$；
$2849.392 - 0.127 - 2846.512 = \qty{2.753}{\milli\gal}$。

\textbf{3.} 各点距基点的高度分别为 $-27.1$、$-30.6$ 和
\qty{-18.1}{m}；加上 $0.1967\,\Delta h$ 后得布格异常
$\num{-0.817}$、$\num{-2.122}$ 和 \qty{-0.808}{\milli\gal}。测点低于基点，且离地心更近，
故其原始读数偏高；
校正去除了这一影响，
在花岗岩上方留下一个低值。

\textbf{4.} 多出的 $0.098 - 0.040 = \qty{0.058}{\milli\gal.h^{-1}}$，
环线上的 \qty{0.12}{\milli\gal} 为固体地球潮汐。两小时内
它近似线性变化，故过两基点读数的直线
可与\omterm{def:b3:field-lab-methods-3:gravimeter}{仪器漂移}一并扣除。

\textbf{5.} 误差分别为 $\num{0.008}$、$\num{0.009}$ 和
\qty{0.017}{\milli\gal}。各高度取整到 \qty{0.1}{m} 时，可使
至多 \qty{0.1}{m} 的高差误差，相当于
\qty{0.020}{\milli\gal}：取整解释了误差，
而直线去除后残留的潮汐小于 \qty{0.001}{\milli\gal}。

\textbf{6.} $0.01/0.1967 = \qty{5.1}{cm}$。手册中的 \qty{3.3}{cm} 则是
$0.01/0.3086$：自由空气异常与布格校正所取的密度无关，
因此高度也须相应取值。

\textbf{7.} $x_{1/2} = (3.952 - 3.048)/2 = \qty{0.452}{km}$，且 $z = 1.305
\times452 = \qty{590}{m}$，即地面以下的深度。

\textbf{8.} $\Delta m = \Delta g\,z^2/G = \num{-2.293e-5}\times590^2/
\num{6.674e-11} = \qty{-1.20e11}{kg}$。取 $\Delta\rho =
\qty{-130}{kg.m^{-3}}$，则其体积为 \qty{9.2e8}{m^3}，半径为
\qty{603}{m}，大于其中心深度，则球体将伸出地面。
不存在具有花岗岩密度反差的球体与之相符；
该体必须接近地表。

\textbf{9.} 花岗岩从地表延伸到很大深度。其浅部
靠近测点，使异常峰变窄；而规则只知宽度，
给出一个位于中等深度的单一球体；深部则
加上球体所缺的宽缓翼部（见剖面图），
其中所含的质量未被计入。

\textbf{10.} $2\pi G\Delta\rho = 2\pi\times\num{6.674e-11}\times(-130) =
\qty{-5.45e-8}{s^{-2}}$。取 $z_1 = 0$，$z_2 = \qty{1645}{m}$，$a =
\qty{461}{m}$：括号内为 $1645 + 461 - 1708.4 = \qty{397.6}{m}$，
$\Delta g = \qty{-2.17}{\milli\gal}$。若深度不设限，括号趋于 $a$，
则 $\Delta g = \qty{-2.51}{\milli\gal}$。

\textbf{11.} 括号须为 $2.29/\num{5.45e-3} = \qty{420}{m}$，
约 \qty{2.6}{km} 处的底面即可满足。任一介于约
\qty{1.6}{km} 与无限深之间的底面，只使峰值改变零点几毫伽，
而略宽或密度稍大的体给出相同峰值，只是底面更浅，
因此底面无法被分辨。

\textbf{12.} 取 $z_1 = \qty{6}{m}$：$1639 + 461.0 - 1708.4 = 391.6$，且
$\Delta g = \qty{-2.14}{\milli\gal}$，比顶面位于地表时少 \qty{0.03}{\milli\gal}，
这小于测量中的地形与密度误差：重力无法判断
花岗岩是否在冲积层之下到达地表；
电测深则可以。

\textbf{13.} $K = \pi(20^2 - 1^2)/(2\times1) = \qty{626.7}{m}$，且 $\Delta V =
\rho_a I/K = 821\times0.050/626.7 = \qty{65.5}{mV}$。

\textbf{14.} $V = 2\times1/\sqrt{17.1^2 - 10.7^2} = 2/13.34 =
\qty{0.150}{m.ns^{-1}}$，$\varepsilon_r = (0.2998/0.150)^2 = 4.0$，为
干砂；管道深度为 $0.150\times10.7/2 = \qty{0.80}{m}$。

\textbf{15.} $0.150\times20.0/2 = \qty{1.50}{m}$。

\textbf{16.} $20/821 = \num{0.0244}$，$30/1309 = \num{0.0229}$，$40/1684 =
\num{0.0238}$：$S \approx \qty{0.0237}{S}$。

\textbf{17.} 干层贡献 $1.5/600 = \qty{0.0025}{S}$，饱和层
贡献 $\qty{0.0212}{S}$，故其厚度为 $0.0212\times222 =
\qty{4.7}{m}$，花岗岩位于 $1.5 + 4.7 = \qty{6.2}{m}$
底面以下。估计值偏高：花岗岩并非无限高阻，
因此在这些间距下 $\rho_a$ 仍低于直线
$s/S$，故 $s/\rho_a$ 高估了 $S$。（该区模型的冲积层厚
\qty{6}{m}。）

\textbf{18.} $\phi = (20/222)^{1/2} = \num{0.30}$。当 $m = 1.3$ 时，$\phi =
0.0901^{1/1.3} = \num{0.16}$；当 $m = 2.8$ 时，$\phi = \num{0.42}$。

\textbf{19.} 将测线减去基线，自西向东依次为：
3.4、4.1、4.8、5.9、7.4、9.2、11.5、14.2、17.2、19.8、21.0、20.1、17.6、
14.5、11.5、9.1、7.2、5.7、4.6、3.8 和 \qty{3.1}{nT}。

\textbf{20.} 未经校正时，峰值为 \qty{28.0}{nT}，位于 \qty{2.20}{km} 处，
其一半 \qty{14.0}{nT} 在东侧从未达到：测线止于
\qty{15.2}{nT}，因为行走期间场强
上升了 \qty{12}{nT}。

\textbf{21.} 峰值为 \qty{21.0}{nT}，位于 \qty{2.20}{km} 处；其一半为
\qty{10.5}{nT}。西侧：\num{1.95} 处的 \num{9.2} 与
\qty{2.00}{km} 处的 \num{11.5} 之间，位于 $1.95 + 0.05\times1.3/2.3 = \qty{1.978}{km}$。东侧：
介于 \num{2.40} 处的 \num{11.5} 与 \qty{2.45}{km} 处的 \num{9.1} 之间，位于
$2.40 + 0.05\times1.0/2.4 = \qty{2.421}{km}$。半宽为 $(2.421
- 1.978)/2 = \qty{0.221}{km}$。

\textbf{22.} 在倾角为 \ang{65} 的场中，对称异常意味着只有
竖直磁化起作用：岩体沿磁子午线走向，
即南–北向。依据
\cref{prop:b3:field-lab-methods-3:dyke}的规则，假定一个薄的、竖直的、
二维岩墙，向下延伸很远，由感应磁化，其顶端
起着磁荷线的作用。

\textbf{23.} $\kappa t = 2\pi z\,\Delta T_\mathrm{max}/(F\sin^2 I) = 2\pi\times
221\times21.0/(48000\times0.8214) = \qty{0.74}{m}$。取 $\kappa =
0.03$，$t = \qty{25}{m}$。只有乘积是确定的：厚度为两倍而磁化率减半的
岩墙给出相同的异常。

\textbf{24.} 异常在峰值一半处宽 \qty{443}{m}，测点相距 \qty{50}{m}，
跨越异常共有九个读数，是规则要求的
最低数的三倍。

\textbf{25.} 岩墙顶部位于
\textbf{\qty[text-series-to-math=true]{220}{m}}以下。此处地质模型
将不整合面置于 \qty{222}{m} 深处：岩墙
上升至覆盖层底部，并在该处被截断，
其年代早于掩埋它的三叠纪砂岩。
\end{solution}
